\documentclass[aspectratio=169,10pt]{beamer}

% Shared Beamer setup for Fall 2026 course slides.
\mode<presentation>{
  \usetheme{Madrid}
  \usecolortheme{default}
}

\definecolor{CalPolyGreen}{HTML}{154734}
\definecolor{CalPolyGold}{HTML}{BD8B13}
\definecolor{DeckInk}{HTML}{1F2933}
\definecolor{DeckMuted}{HTML}{5F6B7A}

\setbeamercolor{structure}{fg=CalPolyGreen}
\setbeamercolor{palette primary}{bg=CalPolyGreen,fg=white}
\setbeamercolor{palette secondary}{bg=DeckInk,fg=white}
\setbeamercolor{palette tertiary}{bg=CalPolyGold,fg=DeckInk}
\setbeamercolor{frametitle}{bg=CalPolyGreen,fg=white}
\setbeamercolor{title}{fg=CalPolyGreen}
\setbeamercolor{normal text}{fg=DeckInk,bg=white}
\setbeamercolor{block title}{bg=CalPolyGreen,fg=white}
\setbeamercolor{block body}{bg=black!3,fg=DeckInk}

\setbeamertemplate{navigation symbols}{}
\setbeamertemplate{footline}[frame number]
\setbeamertemplate{itemize item}{\color{CalPolyGold}$\blacktriangleright$}
\setbeamertemplate{itemize subitem}{\color{CalPolyGreen}$\bullet$}

\newcommand{\CourseNumber}{Course}
\newcommand{\CourseTitle}{Course Title}
\newcommand{\LectureNumber}{0}
\newcommand{\LectureTitle}{Lecture Title}
\newcommand{\LectureDate}{Date}
\newcommand{\CourseUnit}{Unit}

\newcommand{\coursenumber}[1]{\renewcommand{\CourseNumber}{#1}}
\newcommand{\coursetitle}[1]{\renewcommand{\CourseTitle}{#1}}
\newcommand{\lecturenumber}[1]{\renewcommand{\LectureNumber}{#1}}
\newcommand{\lecturetitle}[1]{\renewcommand{\LectureTitle}{#1}}
\newcommand{\lecturedate}[1]{\renewcommand{\LectureDate}{#1}}
\newcommand{\courseunit}[1]{\renewcommand{\CourseUnit}{#1}}

\author{Kirk Duran}
\institute{California Polytechnic State University, San Luis Obispo}
\date{\LectureDate}

\newcommand{\makedecktitle}{
  \begin{frame}[plain]
    \vspace{0.25in}
    {\usebeamerfont{title}\color{CalPolyGreen}\LectureTitle\par}
    \vspace{0.18in}
    {\large \CourseNumber{}: \CourseTitle\par}
    \vspace{0.08in}
    {\large Lecture \LectureNumber{} -- \LectureDate\par}
    \vfill
    {\small Kirk Duran \textbar{} Cal Poly, San Luis Obispo\par}
  \end{frame}
}

\newcommand{\placeholdernote}{
  \begin{block}{Placeholder}
    Replace these bullets with the final lecture material, examples, and in-class activities.
  \end{block}
}

\AtBeginSection[]{
  \begin{frame}{Roadmap}
    \tableofcontents[currentsection]
  \end{frame}
}


\graphicspath{{../assets/lecture-06/}}

% If an image file is not yet available, the slide displays a labeled
% placeholder using the intended filename. Place the matching file in
% ../assets/lecture-06/ to replace the placeholder automatically.
\newcommand{\lectureimage}[2][1.75in]{%
  \IfFileExists{../assets/lecture-06/#2}{%
    \includegraphics[width=\linewidth,height=#1,keepaspectratio]{#2}%
  }{%
    \fbox{%
      \parbox[c][#1][c]{0.94\linewidth}{%
        \centering
        \small\textbf{Image Placeholder}\\[0.08in]
        \scriptsize\texttt{\detokenize{#2}}
      }%
    }%
  }%
}

\setbeamersize{text margin left=0.42in,text margin right=0.42in}

\coursenumber{CS 1001}
\coursetitle{Introduction to Computing}
\lecturenumber{6}
\lecturetitle{Boolean Logic and Comparisons}
\lecturedate{Fall 2026}
\courseunit{1. Computing and Program Execution}

\begin{document}

\makedecktitle

\begin{frame}{Today}
  \begin{itemize}
    \item How Python represents logical truth using the \texttt{bool} type.
    \item How comparison expressions inspect values and produce \texttt{True} or \texttt{False}.
    \item Why assignment with \texttt{=} and equality comparison with \texttt{==} have different semantics.
    \item How \texttt{and}, \texttt{or}, and \texttt{not} combine or negate Boolean claims.
    \item How truth tables, precedence, and parentheses help analyze compound Boolean expressions.
  \end{itemize}

  \begin{keyidea}
    This lecture ends with one capability: Python can compute whether a claim is true or false.
  \end{keyidea}
\end{frame}

\sectionframe{Part 1: Boolean Values}

\begin{frame}{Boolean Logic Begins with Claims}
  \begin{columns}[T,onlytextwidth]
    \begin{column}{0.54\textwidth}
      \begin{itemize}
        \item Previous lectures established that expressions evaluate to values.
        \item Programs also need to represent claims about those values.
        \item A logical claim has one of two outcomes: \texttt{True} or \texttt{False}.
      \end{itemize}

      \begin{block}{Opening expression}
        \centering
        \Large\texttt{8 > 3}
      \end{block}

      \begin{exampleblockcp}
        The expression evaluates to the Boolean value \texttt{True}.
      \end{exampleblockcp}
    \end{column}

    \begin{column}{0.40\textwidth}
      % Search direction: numeric comparison flowing to a True result.
      \lectureimage[2.25in]{comparison-produces-boolean.png}
    \end{column}
  \end{columns}
\end{frame}

\begin{frame}{The \texttt{bool} Type Represents Two Logical Values}
  \begin{cpdefinition}
    A \textbf{Boolean value} represents one of two logical states: \texttt{True} or \texttt{False}.
  \end{cpdefinition}

  \begin{columns}[T,onlytextwidth]
    \begin{column}{0.47\textwidth}
      \begin{block}{Boolean literals}
        \centering
        \Large\texttt{True}\qquad\Large\texttt{False}
      \end{block}
    \end{column}

    \begin{column}{0.47\textwidth}
      \begin{block}{Boolean variables}
        \texttt{is\_raining = True}\\
        \texttt{is\_sunny = False}\\
        \texttt{has\_passed = True}
      \end{block}
    \end{column}
  \end{columns}

  \begin{itemize}
    \item Boolean literals are capitalized in Python.
    \item Their Python type is \texttt{bool}.
    \item Boolean values can be stored, printed, assigned, and combined like other values.
  \end{itemize}
\end{frame}

\begin{frame}{Boolean Values Can Be Stored in Variables}
  \begin{columns}[T,onlytextwidth]
    \begin{column}{0.52\textwidth}
      \begin{block}{Example state}
        \texttt{has\_id = True}\\
        \texttt{has\_ticket = False}\\
        \texttt{is\_late = True}
      \end{block}

      \begin{itemize}
        \item The variable name refers to a Boolean value.
        \item Evaluating the name retrieves its current Boolean value.
        \item This is the same variable model introduced in Lecture 5.
      \end{itemize}
    \end{column}

    \begin{column}{0.42\textwidth}
      % Search direction: variable label has_id pointing to True object.
      \lectureimage[2.25in]{boolean-variable-binding.png}
    \end{column}
  \end{columns}
\end{frame}

\begin{frame}{Booleans and Integers in Python}
  \begin{columns}[T,onlytextwidth]
    \begin{column}{0.52\textwidth}
      \begin{itemize}
        \item In Python, \texttt{bool} is a subtype of \texttt{int}.
        \item Numerically, \texttt{True} behaves like \texttt{1} and \texttt{False} behaves like \texttt{0}.
        \item This implementation relationship is useful context, but the primary abstraction is logical truth.
      \end{itemize}

      \begin{block}{Conceptual binary representation}
        \texttt{True} $\longrightarrow$ \texttt{1}\\
        \texttt{False} $\longrightarrow$ \texttt{0}
      \end{block}
    \end{column}

    \begin{column}{0.42\textwidth}
      \begin{block}{Low-level representation idea}
        \texttt{True} $\longrightarrow$ \texttt{00000001}\\
        \texttt{False} $\longrightarrow$ \texttt{00000000}
      \end{block}

      \begin{alertblock}{Abstraction boundary}
        Python objects have implementation details beyond a single raw bit. Use \texttt{True}/\texttt{False} as the programming model.
      \end{alertblock}
    \end{column}
  \end{columns}
\end{frame}

\sectionframe{Part 2: Comparison Expressions}

\begin{frame}{From Values to Claims About Values}
  \begin{columns}[T,onlytextwidth]
    \begin{column}{0.47\textwidth}
      \begin{block}{A stored value}
        \centering
        \Large\texttt{temperature = 72}
      \end{block}

      \begin{itemize}
        \item The variable \texttt{temperature} refers to the value \texttt{72}.
      \end{itemize}
    \end{column}

    \begin{column}{0.47\textwidth}
      \begin{block}{A claim about the value}
        \centering
        \Large\texttt{temperature > 70}
      \end{block}

      \begin{itemize}
        \item The expression asks whether the current value is greater than \texttt{70}.
        \item The result is \texttt{True}.
      \end{itemize}
    \end{column}
  \end{columns}

  \begin{cpdefinition}
    A \textbf{comparison expression} compares values and evaluates to a Boolean value.
  \end{cpdefinition}
\end{frame}

\begin{frame}{Comparison Expressions Produce Boolean Values}
  \begin{columns}[T,onlytextwidth]
    \begin{column}{0.50\textwidth}
      \begin{block}{Examples}
        \texttt{8 > 3}\hfill$\longrightarrow$\hfill\texttt{True}\\[0.08in]
        \texttt{8 == 3}\hfill$\longrightarrow$\hfill\texttt{False}\\[0.08in]
        \texttt{8 != 3}\hfill$\longrightarrow$\hfill\texttt{True}
      \end{block}
    \end{column}

    \begin{column}{0.44\textwidth}
      \begin{block}{Expression model}
        \centering
        \texttt{8 > 3}\\[0.08in]
        $\downarrow$ evaluate\\[0.08in]
        \texttt{True}
      \end{block}
    \end{column}
  \end{columns}

  \begin{keyidea}
    A comparison is an expression: it computes a value rather than performing a control-flow action.
  \end{keyidea}
\end{frame}

\begin{frame}{Python Provides Six Fundamental Comparison Operators}
  \begin{center}
    \renewcommand{\arraystretch}{1.35}
    \begin{tabular}{c|l@{\qquad}c|l}
      \textbf{Operator} & \textbf{Meaning} & \textbf{Operator} & \textbf{Meaning} \\\hline
      \texttt{<} & less than & \texttt{<=} & less than or equal to \\
      \texttt{>} & greater than & \texttt{>=} & greater than or equal to \\
      \texttt{==} & equal to & \texttt{!=} & not equal to
    \end{tabular}
  \end{center}

  \vspace{0.16in}

  \begin{block}{Examples}
    \texttt{5 > 3} $\rightarrow$ \texttt{True}\qquad
    \texttt{5 < 3} $\rightarrow$ \texttt{False}\qquad
    \texttt{5 == 5} $\rightarrow$ \texttt{True}\\[0.08in]
    \texttt{5 != 3} $\rightarrow$ \texttt{True}\qquad
    \texttt{5 >= 5} $\rightarrow$ \texttt{True}\qquad
    \texttt{5 <= 3} $\rightarrow$ \texttt{False}
  \end{block}
\end{frame}

\begin{frame}{Comparisons Can Inspect Values Stored in Variables}
  \begin{columns}[T,onlytextwidth]
    \begin{column}{0.52\textwidth}
      \begin{block}{Example}
        \texttt{temperature = 72}\\[0.08in]
        \texttt{temperature > 70}
      \end{block}

      \begin{enumerate}
        \item Retrieve the current value of \texttt{temperature}: \texttt{72}.
        \item Compare \texttt{72 > 70}.
        \item Produce the Boolean value \texttt{True}.
      \end{enumerate}
    \end{column}

    \begin{column}{0.42\textwidth}
      % Search direction: substitution/evaluation steps for temperature > 70.
      \lectureimage[2.35in]{comparison-variable-evaluation.png}
    \end{column}
  \end{columns}
\end{frame}

\begin{frame}{Comparison Expressions Do Not Change the Compared Values}
  \begin{block}{Sequence}
    \texttt{temperature = 72}\\
    \texttt{temperature > 70}
  \end{block}

  \begin{columns}[T,onlytextwidth]
    \begin{column}{0.47\textwidth}
      \begin{block}{Comparison result}
        \centering
        \Large\texttt{True}
      \end{block}
    \end{column}

    \begin{column}{0.47\textwidth}
      \begin{block}{Current program state}
        \centering
        \Large\texttt{temperature} $\longrightarrow$ \texttt{72}
      \end{block}
    \end{column}
  \end{columns}

  \begin{keyidea}
    Comparison expressions inspect values and produce Boolean results; they do not themselves reassign variables.
  \end{keyidea}
\end{frame}

\sectionframe{Part 3: Equality Versus Assignment}

\begin{frame}{Assignment and Equality Comparison Have Different Semantics}
  \begin{columns}[T,onlytextwidth]
    \begin{column}{0.47\textwidth}
      \begin{block}{Assignment}
        \centering
        \Large\texttt{x = 5}
      \end{block}

      \begin{itemize}
        \item Evaluates the right-hand side.
        \item Associates \texttt{x} with the resulting value.
        \item Changes program state.
      \end{itemize}
    \end{column}

    \begin{column}{0.47\textwidth}
      \begin{block}{Equality comparison}
        \centering
        \Large\texttt{x == 5}
      \end{block}

      \begin{itemize}
        \item Compares two values.
        \item Produces \texttt{True} or \texttt{False}.
        \item Does not reassign \texttt{x}.
      \end{itemize}
    \end{column}
  \end{columns}
\end{frame}

\begin{frame}{\texttt{=} Changes Program State}
  \begin{columns}[T,onlytextwidth]
    \begin{column}{0.52\textwidth}
      \begin{block}{Assignment}
        \centering
        \Large\texttt{x = 5}
      \end{block}

      \begin{enumerate}
        \item Evaluate the right-hand expression \texttt{5}.
        \item Associate the name \texttt{x} with that value.
      \end{enumerate}
    \end{column}

    \begin{column}{0.42\textwidth}
      % Search direction: assignment changes binding to x -> 5.
      \lectureimage[2.20in]{assignment-x-to-five.png}
    \end{column}
  \end{columns}

  \begin{keyidea}
    Read \texttt{=} as an assignment operation, not as a Boolean test of equality.
  \end{keyidea}
\end{frame}

\begin{frame}{\texttt{==} Produces a Boolean Value}
  \begin{columns}[T,onlytextwidth]
    \begin{column}{0.52\textwidth}
      \begin{block}{Sequence}
        \texttt{x = 5}\\
        \texttt{x == 5}
      \end{block}

      \begin{enumerate}
        \item Evaluate \texttt{x} to retrieve \texttt{5}.
        \item Compare \texttt{5 == 5}.
        \item Produce \texttt{True}.
      \end{enumerate}
    \end{column}

    \begin{column}{0.42\textwidth}
      \begin{block}{Evaluation}
        \centering
        \texttt{x == 5}\\[0.08in]
        $\downarrow$\\[0.08in]
        \texttt{5 == 5}\\[0.08in]
        $\downarrow$\\[0.08in]
        \texttt{True}
      \end{block}
    \end{column}
  \end{columns}
\end{frame}

\begin{frame}{Misconception Check: \texttt{=} Versus \texttt{==}}
  \begin{block}{Code}
    \texttt{x = 10}\\
    \texttt{x == 5}
  \end{block}

  \begin{activity}
    What value does the comparison produce? What value is associated with \texttt{x} afterward?
  \end{activity}

  \begin{columns}[T,onlytextwidth]
    \begin{column}{0.47\textwidth}
      \begin{block}{Comparison result}
        \centering
        \Large\texttt{False}
      \end{block}
    \end{column}

    \begin{column}{0.47\textwidth}
      \begin{block}{Final value of \texttt{x}}
        \centering
        \Large\texttt{10}
      \end{block}
    \end{column}
  \end{columns}

  \begin{keyidea}
    Assignment changes state; comparison produces a value.
  \end{keyidea}
\end{frame}

\sectionframe{Part 4: Boolean Values from Expressions}

\begin{frame}{Boolean Values Can Be Assigned Directly or Computed}
  \begin{columns}[T,onlytextwidth]
    \begin{column}{0.47\textwidth}
      \begin{block}{Direct Boolean assignment}
        \centering
        \texttt{enrolled = True}
      \end{block}
    \end{column}

    \begin{column}{0.47\textwidth}
      \begin{block}{Computed Boolean assignment}
        \texttt{units = 12}\\
        \texttt{enrolled = units >= 1}
      \end{block}
    \end{column}
  \end{columns}

  \vspace{0.12in}

  \begin{block}{Evaluation}
    \centering
    \texttt{units >= 1} $\longrightarrow$ \texttt{12 >= 1} $\longrightarrow$ \texttt{True}
  \end{block}

  \begin{keyidea}
    A comparison expression can produce the Boolean value that assignment then stores.
  \end{keyidea}
\end{frame}

\begin{frame}{Boolean Expressions Use the Same Expression Model}
  \begin{columns}[T,onlytextwidth]
    \begin{column}{0.47\textwidth}
      \begin{block}{Arithmetic expression}
        \texttt{x = 5 + 3}\\[0.08in]
        \texttt{5 + 3} $\longrightarrow$ \texttt{8}
      \end{block}
    \end{column}

    \begin{column}{0.47\textwidth}
      \begin{block}{Boolean expression}
        \texttt{enrolled = units >= 1}\\[0.08in]
        \texttt{units >= 1} $\longrightarrow$ \texttt{True}
      \end{block}
    \end{column}
  \end{columns}

  \begin{enumerate}
    \item Evaluate the expression on the right.
    \item Obtain one resulting value.
    \item Associate the name on the left with that value.
  \end{enumerate}
\end{frame}

\sectionframe{Part 5: Boolean Operators}

\begin{frame}{Boolean Operators Combine or Modify Claims}
  \begin{columns}[T,onlytextwidth]
    \begin{column}{0.54\textwidth}
      \begin{block}{Individual claims}
        \texttt{age >= 18}\\
        \texttt{has\_ticket}
      \end{block}

      \begin{block}{Combined claim}
        \centering
        \Large\texttt{age >= 18 and has\_ticket}
      \end{block}

      \begin{cpdefinition}
        A \textbf{Boolean operator} combines or modifies Boolean values and produces another Boolean result.
      \end{cpdefinition}
    \end{column}

    \begin{column}{0.40\textwidth}
      \begin{block}{Python Boolean operators}
        \centering
        \Large\texttt{and}\\[0.16in]
        \Large\texttt{or}\\[0.16in]
        \Large\texttt{not}
      \end{block}
    \end{column}
  \end{columns}
\end{frame}

\begin{frame}{Boolean AND Requires Both Claims to Be True}
  \begin{columns}[T,onlytextwidth]
    \begin{column}{0.47\textwidth}
      \begin{center}
        \renewcommand{\arraystretch}{1.25}
        \begin{tabular}{c|c|c}
          \textbf{A} & \textbf{B} & \textbf{A and B} \\\hline
          False & False & False \\
          False & True & False \\
          True & False & False \\
          True & True & True
        \end{tabular}
      \end{center}
    \end{column}

    \begin{column}{0.47\textwidth}
      \begin{block}{Example}
        \texttt{has\_id = True}\\
        \texttt{has\_ticket = True}\\[0.08in]
        \texttt{has\_id and has\_ticket}\\[0.08in]
        $\longrightarrow$ \texttt{True}
      \end{block}
    \end{column}
  \end{columns}

  \begin{keyidea}
    \texttt{and} is true only when both component claims are true.
  \end{keyidea}
\end{frame}

\begin{frame}{Boolean OR Requires At Least One True Claim}
  \begin{columns}[T,onlytextwidth]
    \begin{column}{0.47\textwidth}
      \begin{center}
        \renewcommand{\arraystretch}{1.25}
        \begin{tabular}{c|c|c}
          \textbf{A} & \textbf{B} & \textbf{A or B} \\\hline
          False & False & False \\
          False & True & True \\
          True & False & True \\
          True & True & True
        \end{tabular}
      \end{center}
    \end{column}

    \begin{column}{0.47\textwidth}
      \begin{block}{Example}
        \texttt{is\_weekend = False}\\
        \texttt{is\_holiday = True}\\[0.08in]
        \texttt{is\_weekend or is\_holiday}\\[0.08in]
        $\longrightarrow$ \texttt{True}
      \end{block}
    \end{column}
  \end{columns}

  \begin{keyidea}
    Python's logical \texttt{or} is inclusive: it is also true when both component claims are true.
  \end{keyidea}
\end{frame}

\begin{frame}{Boolean NOT Reverses a Truth Value}
  \begin{columns}[T,onlytextwidth]
    \begin{column}{0.40\textwidth}
      \begin{center}
        \renewcommand{\arraystretch}{1.35}
        \begin{tabular}{c|c}
          \textbf{A} & \textbf{not A} \\\hline
          False & True \\
          True & False
        \end{tabular}
      \end{center}
    \end{column}

    \begin{column}{0.54\textwidth}
      \begin{block}{Example}
        \texttt{is\_late = True}\\[0.08in]
        \texttt{not is\_late}\\[0.08in]
        $\longrightarrow$ \texttt{False}
      \end{block}

      \begin{itemize}
        \item \texttt{not} is a unary Boolean operator.
        \item It negates the truth value of one Boolean operand.
      \end{itemize}
    \end{column}
  \end{columns}
\end{frame}

\begin{frame}{Read Boolean Expressions as Logical Claims}
  \begin{center}
    \renewcommand{\arraystretch}{1.35}
    \begin{tabular}{l|l}
      \textbf{Python} & \textbf{Plain-language claim} \\\hline
      \texttt{has\_id and has\_ticket} & has identification and has a ticket \\
      \texttt{is\_weekend or is\_holiday} & it is the weekend or a holiday \\
      \texttt{not is\_late} & it is not late \\
      \texttt{temperature >= 70 and is\_sunny} & at least 70 degrees and sunny
    \end{tabular}
  \end{center}

  \begin{activity}
    Translate: ``The student has completed the prerequisite and is enrolled.''
  \end{activity}
\end{frame}

\sectionframe{Part 6: Truth Tables}

\begin{frame}{Truth Tables Enumerate Every Boolean Input Combination}
  \begin{cpdefinition}
    A \textbf{truth table} lists every possible combination of Boolean inputs and the result of a logical expression for each combination.
  \end{cpdefinition}

  \begin{itemize}
    \item Truth tables provide a systematic method for analyzing Boolean expressions.
    \item They clarify the exact semantics of \texttt{and}, \texttt{or}, and \texttt{not}.
    \item They are useful for reasoning about and debugging compound logical statements.
  \end{itemize}

  \begin{keyidea}
    For two Boolean inputs, there are four possible input combinations.
  \end{keyidea}
\end{frame}

\begin{frame}{Combined Truth Table for \texttt{and} and \texttt{or}}
  \begin{center}
    \renewcommand{\arraystretch}{1.45}
    \begin{tabular}{c|c|c|c}
      \textbf{A} & \textbf{B} & \textbf{A and B} & \textbf{A or B} \\\hline
      True & True & True & True \\
      True & False & False & True \\
      False & True & False & True \\
      False & False & False & False
    \end{tabular}
  \end{center}

  \begin{keyidea}
    The table encodes the complete behavior of each binary Boolean operator for Boolean operands.
  \end{keyidea}
\end{frame}

\begin{frame}{Truth Tables Can Analyze Compound Expressions}
  \begin{block}{Expression}
    \centering
    \Large\texttt{A and not B}
  \end{block}

  \begin{center}
    \renewcommand{\arraystretch}{1.35}
    \begin{tabular}{c|c|c|c}
      \textbf{A} & \textbf{B} & \textbf{not B} & \textbf{A and not B} \\\hline
      True & True & False & False \\
      True & False & True & True \\
      False & True & False & False \\
      False & False & True & False
    \end{tabular}
  \end{center}
\end{frame}

\begin{frame}{Use Truth Tables to Reason Systematically}
  \begin{columns}[T,onlytextwidth]
    \begin{column}{0.30\textwidth}
      \begin{block}{\texttt{A and B}}
        Ask whether both claims are true.
      \end{block}
    \end{column}

    \begin{column}{0.30\textwidth}
      \begin{block}{\texttt{A or B}}
        Ask whether at least one claim is true.
      \end{block}
    \end{column}

    \begin{column}{0.30\textwidth}
      \begin{block}{\texttt{not A}}
        Reverse the truth value of \texttt{A}.
      \end{block}
    \end{column}
  \end{columns}

  \begin{activity}
    Without running Python, determine the result of \texttt{True and not False} and justify each step.
  \end{activity}
\end{frame}

\sectionframe{Part 7: Compound Boolean Expressions}

\begin{frame}{Compound Boolean Expressions Combine Multiple Claims}
  \begin{cpdefinition}
    A \textbf{compound Boolean expression} combines two or more Boolean expressions using Boolean operators and evaluates to one Boolean value.
  \end{cpdefinition}

  \begin{block}{Example}
    \texttt{age = 19}\\
    \texttt{has\_ticket = True}\\[0.08in]
    \texttt{age >= 18 and has\_ticket}
  \end{block}

  \begin{block}{Evaluation}
    \centering
    \texttt{19 >= 18 and True} $\longrightarrow$ \texttt{True and True} $\longrightarrow$ \texttt{True}
  \end{block}
\end{frame}

\begin{frame}{Evaluate Compound Expressions in Stages}
  \begin{columns}[T,onlytextwidth]
    \begin{column}{0.52\textwidth}
      \begin{block}{Example}
        \texttt{temperature = 72}\\
        \texttt{is\_sunny = True}\\[0.08in]
        \texttt{temperature >= 70 and is\_sunny}
      \end{block}

      \begin{enumerate}
        \item Substitute current variable values.
        \item Evaluate comparisons.
        \item Evaluate Boolean operators.
      \end{enumerate}
    \end{column}

    \begin{column}{0.42\textwidth}
      \begin{block}{Stages}
        \texttt{72 >= 70 and True}\\
        $\downarrow$\\
        \texttt{True and True}\\
        $\downarrow$\\
        \texttt{True}
      \end{block}
    \end{column}
  \end{columns}
\end{frame}

\begin{frame}{Example: Assigning the Result of a Compound Expression}
  \begin{block}{Statement}
    \centering
    \Large\texttt{x = 3 > 4 and 5 < 6}
  \end{block}

  \begin{enumerate}
    \item \texttt{3 > 4} $\longrightarrow$ \texttt{False}
    \item \texttt{5 < 6} $\longrightarrow$ \texttt{True}
    \item \texttt{False and True} $\longrightarrow$ \texttt{False}
    \item Assignment associates \texttt{x} with \texttt{False}.
  \end{enumerate}

  \begin{keyidea}
    The right-hand side is fully evaluated before its Boolean result is assigned.
  \end{keyidea}
\end{frame}

\begin{frame}[shrink=6]{Practice: Evaluate Compound Logical Statements}
  \begin{block}{Expressions from the prior lecture deck}
    \texttt{x = 3 > 4 and 5 < 6}\\
    \texttt{y = 4 == 4 or 7 != 7}\\
    \texttt{z = not (10 <= 15)}\\
    \texttt{a = (8 > 2 and 3 < 5) or (6 == 6)}\\
    \texttt{b = not (5 > 2) and (3 >= 3)}\\
    \texttt{c = 7 != 7 or (2 + 2 == 4 and 9 < 10)}\\
    \texttt{d = (3 < 1) and not (5 == 5)}
  \end{block}

  \begin{activity}
    Evaluate each expression hierarchically: arithmetic $\rightarrow$ comparisons $\rightarrow$ Boolean operations $\rightarrow$ final Boolean value.
  \end{activity}
\end{frame}

\sectionframe{Part 8: Operator Precedence and Parentheses}

\begin{frame}{Operator Precedence Determines Evaluation Order}
  \begin{block}{Practical precedence for today's operators}
    \begin{enumerate}
      \item Arithmetic
      \item Comparisons
      \item \texttt{not}
      \item \texttt{and}
      \item \texttt{or}
    \end{enumerate}
  \end{block}

  \begin{block}{Example}
    \texttt{age + 1 >= 18 and has\_ticket}
  \end{block}

  \begin{exampleblockcp}
    Conceptually: \texttt{((age + 1) >= 18) and has\_ticket}
  \end{exampleblockcp}
\end{frame}

\begin{frame}{Parentheses Make Logical Structure Explicit}
  \begin{block}{Expression}
    \centering
    \Large\texttt{(age >= 18 and has\_ticket) or has\_parent\_permission}
  \end{block}

  \begin{columns}[T,onlytextwidth]
    \begin{column}{0.55\textwidth}
      \begin{itemize}
        \item Parentheses can override normal precedence.
        \item Even when they do not change the result, they can make intended grouping easier to read.
        \item Use parentheses when they clarify the logical structure of a compound claim.
      \end{itemize}
    \end{column}

    \begin{column}{0.39\textwidth}
      % Search direction: expression tree with grouped conjunction feeding an OR node.
      \lectureimage[2.25in]{boolean-parentheses-expression-tree.png}
    \end{column}
  \end{columns}
\end{frame}

\sectionframe{Part 9: Translating and Evaluating Claims}

\begin{frame}{Translate Plain-Language Claims into Python}
  \begin{center}
    \renewcommand{\arraystretch}{1.45}
    \begin{tabular}{p{0.49\textwidth}|p{0.43\textwidth}}
      \textbf{Claim} & \textbf{Python expression} \\\hline
      The person is at least 18 and has a ticket. & \texttt{age >= 18 and has\_ticket} \\
      Today is a weekend or a holiday. & \texttt{is\_weekend or is\_holiday} \\
      The student is not late. & \texttt{not is\_late} \\
      The temperature is at least 70 and it is sunny. & \texttt{temperature >= 70 and is\_sunny}
    \end{tabular}
  \end{center}
\end{frame}

\begin{frame}{Arithmetic Expressions Can Appear Inside Comparisons}
  \begin{block}{Example}
    \texttt{score = 84}\\
    \texttt{bonus = 7}\\[0.08in]
    \texttt{score + bonus >= 90}
  \end{block}

  \begin{block}{Evaluation}
    \centering
    \texttt{84 + 7 >= 90} $\longrightarrow$ \texttt{91 >= 90} $\longrightarrow$ \texttt{True}
  \end{block}

  \begin{keyidea}
    Larger expressions are built from smaller expressions; arithmetic may produce values that comparisons then inspect.
  \end{keyidea}
\end{frame}

\begin{frame}{A Complex Boolean Expression Still Produces One Value}
  \begin{block}{State}
    \texttt{age = 19}\\
    \texttt{has\_ticket = False}\\
    \texttt{has\_parent\_permission = True}
  \end{block}

  \begin{block}{Expression}
    \texttt{(age >= 18 and has\_ticket) or has\_parent\_permission}
  \end{block}

  \begin{block}{Evaluation}
    \centering
    \texttt{(True and False) or True} $\longrightarrow$ \texttt{False or True} $\longrightarrow$ \texttt{True}
  \end{block}
\end{frame}

\sectionframe{Part 10: Misconceptions to Correct Explicitly}

\begin{frame}{Misconception: \texttt{=} and \texttt{==} Both Mean ``Equals''}
  \begin{columns}[T,onlytextwidth]
    \begin{column}{0.47\textwidth}
      \begin{block}{\texttt{x = 5}}
        Assignment associates \texttt{x} with \texttt{5}; it changes program state.
      \end{block}
    \end{column}

    \begin{column}{0.47\textwidth}
      \begin{block}{\texttt{x == 5}}
        Equality comparison asks whether two values are equal; it produces a Boolean value.
      \end{block}
    \end{column}
  \end{columns}

  \begin{keyidea}
    Similar symbols do not imply similar semantics.
  \end{keyidea}
\end{frame}

\begin{frame}{Misconception: A Comparison Changes the Variable}
  \begin{block}{Code}
    \texttt{x = 5}\\
    \texttt{x > 3}
  \end{block}

  \begin{columns}[T,onlytextwidth]
    \begin{column}{0.47\textwidth}
      \begin{block}{Expression result}
        \centering
        \Large\texttt{True}
      \end{block}
    \end{column}

    \begin{column}{0.47\textwidth}
      \begin{block}{Value still associated with \texttt{x}}
        \centering
        \Large\texttt{5}
      \end{block}
    \end{column}
  \end{columns}
\end{frame}

\begin{frame}{Misconception: \texttt{and} Means Either Claim May Be True}
  \begin{block}{Correct rule}
    \centering
    \Large\texttt{A and B} is \texttt{True} only when \texttt{A} and \texttt{B} are both \texttt{True}.
  \end{block}

  \begin{center}
    \texttt{True and True} $\longrightarrow$ \texttt{True}
  \end{center}

  \begin{keyidea}
    Every other Boolean input combination for \texttt{and} produces \texttt{False}.
  \end{keyidea}
\end{frame}

\begin{frame}{Misconception: \texttt{or} Requires Exactly One True Claim}
  \begin{block}{Correct rule}
    \centering
    \Large\texttt{A or B} is \texttt{True} when at least one operand is \texttt{True}.
  \end{block}

  \begin{center}
    \texttt{True or True} $\longrightarrow$ \texttt{True}
  \end{center}

  \begin{keyidea}
    Logical \texttt{or} is inclusive, not exclusive.
  \end{keyidea}
\end{frame}

\begin{frame}{Misconception: Boolean Variables Must Be Assigned Literals}
  \begin{columns}[T,onlytextwidth]
    \begin{column}{0.47\textwidth}
      \begin{block}{Literal value}
        \centering
        \texttt{enrolled = True}
      \end{block}
    \end{column}

    \begin{column}{0.47\textwidth}
      \begin{block}{Computed value}
        \texttt{units = 12}\\
        \texttt{enrolled = units >= 1}
      \end{block}
    \end{column}
  \end{columns}

  \begin{keyidea}
    Variables can store Boolean values whether those values are written directly or computed by expressions.
  \end{keyidea}
\end{frame}

\sectionframe{Part 11: Extensions from the Prior Deck}

\begin{frame}{Exclusive OR Distinguishes ``At Least One'' from ``Exactly One''}
  \begin{columns}[T,onlytextwidth]
    \begin{column}{0.47\textwidth}
      \begin{center}
        \renewcommand{\arraystretch}{1.25}
        \begin{tabular}{c|c|c}
          \textbf{A} & \textbf{B} & \textbf{A XOR B} \\\hline
          False & False & False \\
          False & True & True \\
          True & False & True \\
          True & True & False
        \end{tabular}
      \end{center}
    \end{column}

    \begin{column}{0.47\textwidth}
      \begin{block}{Conceptual distinction}
        \textbf{OR}: at least one claim is true.\\[0.12in]
        \textbf{XOR}: exactly one claim is true.
      \end{block}

      \begin{exampleblockcp}
        In Python, \texttt{\textasciicircum} is the bitwise XOR operator and also operates on Boolean operands, but \texttt{and}, \texttt{or}, and \texttt{not} are the primary logical operators for this lecture.
      \end{exampleblockcp}
    \end{column}
  \end{columns}
\end{frame}

\begin{frame}{George Boole and the Origins of Boolean Algebra}
  \begin{columns}[T,onlytextwidth]
    \begin{column}{0.56\textwidth}
      \begin{itemize}
        \item George Boole (1815--1864) was an English mathematician and logician.
        \item He developed an algebraic system for reasoning about logical propositions.
        \item Boolean algebra became foundational to formal logic, digital circuits, and computing.
        \item Claude Shannon later connected Boolean algebra to electrical switching circuits.
        \item The Python type name \texttt{bool} reflects this logical tradition.
      \end{itemize}
    \end{column}

    \begin{column}{0.38\textwidth}
      % Search direction: public-domain portrait of George Boole.
      \lectureimage[2.75in]{george-boole-portrait.png}
    \end{column}
  \end{columns}
\end{frame}

\sectionframe{Part 12: Lecture Boundary and Synthesis}

\begin{frame}{What Boolean Logic Can Compute}
  \begin{block}{Examples}
    \texttt{temperature > 70}\\[0.06in]
    \texttt{age >= 18 and has\_ticket}\\[0.06in]
    \texttt{is\_weekend or is\_holiday}\\[0.06in]
    \texttt{not is\_late}
  \end{block}

  \begin{keyidea}
    Every expression above evaluates to one Boolean value: \texttt{True} or \texttt{False}.
  \end{keyidea}
\end{frame}

\begin{frame}{The Boundary: Do Not Introduce Control Flow Yet}
  \begin{itemize}
    \item Do not introduce \texttt{if} in the core lecture.
    \item Do not describe Boolean expressions as choosing which code runs.
    \item Do not require short-circuit evaluation yet.
    \item Do not introduce chained comparisons.
    \item Do not introduce membership operators.
  \end{itemize}

  \begin{columns}[T,onlytextwidth]
    \begin{column}{0.47\textwidth}
      \begin{block}{Lecture 6 ends here}
        Python can compute whether a claim is true.
      \end{block}
    \end{column}

    \begin{column}{0.47\textwidth}
      \begin{block}{Next conceptual step}
        Python can use that Boolean result to choose an action.
      \end{block}
    \end{column}
  \end{columns}
\end{frame}

\begin{frame}{Boolean Logic: Complete Mental Model}
  \begin{center}
    \Large
    Values
    $\longrightarrow$
    Variables
    $\longrightarrow$
    Expressions
    $\longrightarrow$
    Comparisons
    $\longrightarrow$
    Boolean values
  \end{center}

  \vspace{0.18in}

  \begin{center}
    \Large
    Boolean values
    $\longrightarrow$
    Boolean operators
    $\longrightarrow$
    Compound Boolean expressions
  \end{center}

  \vspace{0.18in}

  \begin{block}{Representative examples}
    \texttt{temperature = 72}\qquad\texttt{temperature > 70}\\[0.08in]
    \texttt{age = 19}\qquad\texttt{age >= 18 and has\_ticket}\\[0.08in]
    \texttt{units = 12}\qquad\texttt{enrolled = units >= 1}
  \end{block}
\end{frame}

\begin{frame}{Lecture 6 Learning Outcome}
  \begin{itemize}
    \item Identify and store Boolean values.
    \item Use the six comparison operators correctly.
    \item Distinguish assignment from equality comparison.
    \item Evaluate comparison expressions involving literal values and variables.
    \item Use \texttt{and}, \texttt{or}, and \texttt{not} to form compound Boolean expressions.
    \item Interpret truth tables and evaluate compound expressions systematically.
    \item Apply practical operator precedence and use parentheses to communicate structure.
    \item Explain that Boolean expressions compute values without yet invoking control flow.
  \end{itemize}

  \begin{keyidea}
    Outcome: students can evaluate simple and compound Boolean expressions involving literal values and variables.
  \end{keyidea}
\end{frame}

\begin{frame}{Exit Ticket: Evaluate and Explain}
  \begin{block}{State}
    \texttt{age = 17}\\
    \texttt{has\_ticket = True}\\
    \texttt{has\_parent\_permission = True}
  \end{block}

  \begin{activity}
    Evaluate \texttt{(age >= 18 and has\_ticket) or has\_parent\_permission}. Show each intermediate Boolean value and explain why evaluating the expression does not change \texttt{age} or \texttt{has\_ticket}.
  \end{activity}
\end{frame}

\appendix
\sectionframe{Appendix: Deferred Material from the Prior Lecture 6 Deck}

\begin{frame}{Boolean Logic in the CPU}
  \begin{columns}[T,onlytextwidth]
    \begin{column}{0.54\textwidth}
      \begin{itemize}
        \item Digital processors implement logical operations using electronic circuits corresponding to operations such as AND, OR, and NOT.
        \item Intermediate computational results may be stored temporarily in CPU registers.
        \item At the hardware level, logical states are represented using binary electrical states.
      \end{itemize}

      \begin{block}{Conceptual connection}
        Boolean expression $\longrightarrow$ logical operations $\longrightarrow$ binary result $\longrightarrow$ true/false interpretation
      \end{block}
    \end{column}

    \begin{column}{0.40\textwidth}
      % Search direction: simple AND/OR/NOT logic-gate diagram.
      \lectureimage[2.55in]{cpu-boolean-logic-gates.png}
    \end{column}
  \end{columns}
\end{frame}

\begin{frame}{Branching Instructions Use Boolean-Like Results}
  \begin{itemize}
    \item Machine instructions can alter execution order based on condition results.
    \item At a high level, a true condition may select one execution path and a false condition another.
    \item This is the conceptual bridge from Boolean evaluation to control flow.
  \end{itemize}

  \begin{alertblock}{Placement}
    This mechanism belongs with the next control-flow unit. The core Lecture 6 stops before introducing \texttt{if} or program branching.
  \end{alertblock}
\end{frame}

\begin{frame}{Short-Circuit Evaluation}
  \begin{itemize}
    \item Python's \texttt{and} and \texttt{or} operators use short-circuit evaluation.
    \item For \texttt{and}, evaluation can stop once an operand determines that the complete expression cannot be true.
    \item For \texttt{or}, evaluation can stop once an operand determines that the complete expression is already true.
  \end{itemize}

  \begin{block}{Examples from the prior deck}
    \texttt{3 < 2 and 5 < 6 and 7 == 7 and 10 >= 10 and 12 != 13}\\[0.10in]
    \texttt{3 > 5 or 6 < 4 or 7 == 7 or 10 < 9 or 12 != 12}
  \end{block}

  \begin{alertblock}{Placement}
    Defer this behavior until students need to reason about evaluation side effects, function calls, or control-flow conditions.
  \end{alertblock}
\end{frame}

\begin{frame}{De Morgan's Laws}
  \begin{columns}[T,onlytextwidth]
    \begin{column}{0.47\textwidth}
      \begin{block}{First law}
        \texttt{not (A and B)}\\[0.08in]
        is logically equivalent to\\[0.08in]
        \texttt{not A or not B}
      \end{block}
    \end{column}

    \begin{column}{0.47\textwidth}
      \begin{block}{Second law}
        \texttt{not (A or B)}\\[0.08in]
        is logically equivalent to\\[0.08in]
        \texttt{not A and not B}
      \end{block}
    \end{column}
  \end{columns}

  \begin{itemize}
    \item De Morgan's Laws transform negated compound Boolean expressions into logically equivalent forms.
    \item They are useful for simplifying and reasoning about more complex logical conditions.
  \end{itemize}
\end{frame}

\begin{frame}{Image Placeholder Index}
  \scriptsize
  \begin{columns}[T,onlytextwidth]
    \begin{column}{0.49\textwidth}
      \begin{itemize}
        \item \texttt{comparison-produces-boolean.png}
        \item \texttt{boolean-variable-binding.png}
        \item \texttt{comparison-variable-evaluation.png}
        \item \texttt{assignment-x-to-five.png}
        \item \texttt{boolean-parentheses-expression-tree.png}
      \end{itemize}
    \end{column}

    \begin{column}{0.49\textwidth}
      \begin{itemize}
        \item \texttt{george-boole-portrait.png}
        \item \texttt{cpu-boolean-logic-gates.png}
      \end{itemize}
    \end{column}
  \end{columns}

  \begin{exampleblockcp}
    Save each image under \texttt{../assets/lecture-06/}. The labeled placeholder will be replaced automatically when its filename matches.
  \end{exampleblockcp}
\end{frame}

\end{document}
